\section*{NeurIPS Paper Checklist}

\begin{enumerate}

\item {\bf Claims}
    \item[] Question: Do the main claims made in the abstract and introduction accurately reflect the paper's contributions and scope?
    \item[] Answer: \answerYes{}
    \item[] Justification: The abstract and Introduction state the contributions (two-tower recommender, prospective temporal evaluation, repurposing candidate set, interactive platform) that are quantitatively supported by Tables~\ref{tab:cv}--\ref{tab:shortlist} and \S\ref{sec:results}.

\item {\bf Limitations}
    \item[] Question: Does the paper discuss the limitations of the work performed by the authors?
    \item[] Answer: \answerYes{}
    \item[] Justification: A dedicated \S\ref{sec:limitations} discusses auxiliary-head leakage, the single temporal-split run, illustrative (not systematically validated) discovery examples, the BoW text encoder, the implicit-feedback PU formulation, and the druggable-genome inference scope.

\item {\bf Theory assumptions and proofs}
    \item[] Question: For each theoretical result, does the paper provide the full set of assumptions and a complete (and correct) proof?
    \item[] Answer: \answerNA{}
    \item[] Justification: The paper is empirical and contains no theoretical results.

\item {\bf Experimental result reproducibility}
    \item[] Question: Does the paper fully disclose all the information needed to reproduce the main experimental results?
    \item[] Answer: \answerYes{}
    \item[] Justification: Architecture, optimizer, learning rate, batch size, dropout range, validation fraction, splits, baselines, and OTP release versions are reported in \S\ref{sec:methods}; full code and trained weights are released (Availability).

\item {\bf Open access to data and code}
    \item[] Question: Does the paper provide open access to the data and code, with sufficient instructions to faithfully reproduce the main experimental results?
    \item[] Answer: \answerYes{}
    \item[] Justification: Code, preprocessed data, model weights, and ranked predictions are released under MIT at the repository linked in the Availability section; OTP source data is publicly available under CC0.

\item {\bf Experimental setting/details}
    \item[] Question: Does the paper specify all the training and test details necessary to understand the results?
    \item[] Answer: \answerYes{}
    \item[] Justification: Splits ($5{\times}5$ target-disjoint CV; 2022$\to$2025 temporal split with leakage filter), filtering criteria, optimizer (Adam, LR $8{\times}10^{-3}$, batch 1{,}024), early stopping, and baseline configurations are specified in \S\ref{sec:methods} and \S\ref{sec:eval}.

\item {\bf Experiment statistical significance}
    \item[] Question: Does the paper report error bars suitably and correctly defined or other appropriate information about statistical significance?
    \item[] Answer: \answerYes{}
    \item[] Justification: Cross-validation results report mean and standard deviation across 25 target-disjoint folds (Table~\ref{tab:cv}); enrichment statistics are Bonferroni-corrected for the six tests performed (\S\ref{sec:char}). The temporal split is a single run; this limitation is explicitly stated in \S\ref{sec:limitations}.

\item {\bf Experiments compute resources}
    \item[] Question: Does the paper provide sufficient information on the computer resources needed?
    \item[] Answer: \answerYes{}
    \item[] Justification: Each cross-validation fold trains in under 30~min on a single consumer GPU (NVIDIA RTX-class, 16~GB); the full $5{\times}5$ CV completes overnight. Inference over the $\sim$4{,}600 druggable-genome targets $\times$ $\sim$19{,}000 diseases ($\sim$87M pairs) runs in a few hours on the same hardware. CPU-only training is feasible but slower.

\item {\bf Code of ethics}
    \item[] Question: Does the research conform with the NeurIPS Code of Ethics?
    \item[] Answer: \answerYes{}
    \item[] Justification: The work uses only public, properly licensed biomedical resources (Open Targets, EFO, Reactome, UniProt, gnomAD); no human-subject data or personally identifiable information is involved.

\item {\bf Broader impacts}
    \item[] Question: Does the paper discuss both potential positive societal impacts and negative societal impacts?
    \item[] Answer: \answerYes{}
    \item[] Justification: Positive impacts (orphan-disease candidate generation, reduction of streetlight bias) and risks (mis-prioritization in resource-limited settings, automation bias) are discussed in \S\ref{sec:limitations} and the Conclusion, which explicitly frames OTRec as a screening tool that does not replace experimental validation.

\item {\bf Safeguards}
    \item[] Question: Does the paper describe safeguards for responsible release of high-misuse-risk data or models?
    \item[] Answer: \answerNA{}
    \item[] Justification: The released model is a target-prioritization recommender over public biomedical data; it is not a generative model, weapons-relevant capability, or content-generation system, and poses no identified high-misuse risk.

\item {\bf Licenses for existing assets}
    \item[] Question: Are the creators or original owners of assets used in the paper properly credited and are the license and terms of use explicitly mentioned and properly respected?
    \item[] Answer: \answerYes{}
    \item[] Justification: Open Targets (CC0), EFO (Apache 2.0), gnomAD (ODbL), Reactome (CC-BY 4.0), UniProt (CC-BY 4.0), TensorFlow/Keras (Apache 2.0), CatBoost (Apache 2.0), and XGBoost (Apache 2.0) are all cited; the released code uses MIT.

\item {\bf New assets}
    \item[] Question: Are new assets introduced in the paper well documented?
    \item[] Answer: \answerYes{}
    \item[] Justification: The released model weights, training scripts, and ranked candidate tables are documented in the repository README (architecture, intended use, training procedure, license, known limitations).

\item {\bf Crowdsourcing and research with human subjects}
    \item[] Question: For crowdsourcing experiments and research with human subjects, does the paper include the full text of instructions given to participants, screenshots, and details about compensation?
    \item[] Answer: \answerNA{}
    \item[] Justification: The work involves no crowdsourcing or human-subject research.

\item {\bf Institutional review board (IRB) approvals or equivalent}
    \item[] Question: Does the paper describe potential risks incurred by study participants and IRB approvals?
    \item[] Answer: \answerNA{}
    \item[] Justification: The work involves no human-subject research and uses only public, de-identified biomedical knowledge bases.

\item {\bf Declaration of LLM usage}
    \item[] Question: Does the paper describe the usage of LLMs if it is an important, original, or non-standard component of the core methods?
    \item[] Answer: \answerNA{}
    \item[] Justification: No LLM is part of the OTRec method; LLMs were used only for incidental writing/editing assistance, which the NeurIPS policy does not require declaring.

\end{enumerate}
